\part{绪论}

\section{系统概览}

\subsection{研究背景与意义}

\subsubsection{气垫船的技术特点与控制需求}

气垫船是一种利用空气动力学原理实现高效行驶的全地形载具。它的核心工作原理是通过底部安装的风扇或高速喷嘴系统，向下持续喷射高压空气，在船体与水面或地面之间形成一个稳定的空气垫层，从而使船体被抬升离开接触面一定高度。这个空气垫不仅有效支撑船体重量，还大幅减少了船体与行进表面之间的接触面积，将传统的滑动摩擦或水体阻力转化为空气与表面的低摩擦效应，从而显著降低运动阻力。

得益于这种独特的设计，气垫船能够获得比传统水上船只更高的航行速度，同时在多种复杂地形上保持优越的通过性。它不仅可以在开阔水域快速航行，还能轻松穿越泥滩、浅滩、湿地、沼泽、冰面、雪地、沙滩甚至相对平坦的陆地，实现水陆两栖运行。这一特性使气垫船摆脱了传统船舶对深水航道的依赖，也克服了轮式或履带车辆对地面平整度的要求，成为真正意义上的全地形运输平台。

在军事领域，气垫船常被用于两栖登陆作战、快速兵力投送、后勤支援及巡逻警戒任务。其高速机动性与地形适应性可在登陆行动中避开潮汐、滩涂障碍与水下防御工事，实现突然、多方向的进攻，显著提升作战效能。在民用领域，气垫船在应急救援、海事搜救、湿地科考、旅游观光及短途客运等方面发挥着不可替代的作用。在洪水、冰灾、沼泽等难以通行的灾害现场，气垫船能够快速抵达并展开救援，运送人员与物资；在生态脆弱的湿地、冰川区域，它因其对地表低破坏的特性，被广泛用于环境监测与科学研究；此外，一些景区也利用气垫船开展观光体验项目，让游客感受在多种地形上"贴地飞行"的独特体验。

从技术发展角度看，气垫船代表了人类在交通工具设计中突破介质限制的重要尝试，融合了船舶与航空器的部分技术理念。然而，其卓越性能的实现，高度依赖于一个响应迅速、稳定可靠、多任务协同的智能控制系统。该系统需实时协调大功率升力风扇、推进舵机、姿态传感器与操控机构，以维持气垫厚度、航行姿态与航向的稳定。因此，开发一套高性能、低功耗、高可靠性的嵌入式控制系统，是提升气垫船整体性能、拓展其应用范围的核心关键技术。

\subsubsection{RISC-V 架构的发展及其在嵌入式领域的优势}

\subsection{国内外研究现状}

\subsubsection{气垫船控制系统发展概况}

\subsubsection{RISC-V在工业控制中的应用现状}

\subsection{本文的主要工作}